\subsection{Ablation: Paired versus Single-Sided Interventions}
\label{sec:paired_ablation}

Does subtracting a matched competitor improve category-direction estimation
beyond subtracting the unedited image? We isolate this choice by fitting four
estimators on identical intervention triples and evaluating them with the
same cosine readout and top-5\% pooling rule
(\cref{sec:presence_evaluation}). This ablation reports raw macro ROC-AUC
and macro average precision (AP).

\paragraph{Controlled comparison.}
We use frozen Qwen2.5-VL-3B-Instruct with the prompt \texttt{Describe.}
and 2,104 accepted intervention records covering 10 object categories and
all 45 unordered category pairs. The saved records are replayed with fixed
donors, base images, placement parameters, and historical selection
decisions. All variants use the same hidden states and the same 5,000
COCO val2017 evaluation images. L18 is the primary layer; L1, L9, L27,
and L35 provide secondary checks. These choices were specified before this
ablation run, after earlier results on the validation population had been
examined.

Let \(H_b,H_c,H_q\) denote the base, target-composite, and
competitor-composite token states. P0, P1, and P2 share the response weights
\(\alpha\) in \cref{eq:response_weights}. Their target contributions are
\(\alpha^\top(H_c-H_q)\), \(\alpha^\top(H_c-H_b)\), and
\(\alpha^\top H_c\), respectively. P3 uses the single-sided difference
with its own weights,
\(\alpha_{c,t}=e_{c,t}/(\sum_s e_{c,s}+10^{-8})\), where
\(e_{c,t}=\tfrac12\|H_{c,t}-H_{b,t}\|_2\).
Thus P0 versus P1 isolates the subtraction reference under fixed weights,
whereas P3 tests whether weights computed without the opposite composite
recover the paired result. P2 removes subtraction from the pooled vector;
its weights still depend on interventions.

Each record contributes to both endpoint categories. Only P0 uses opposite
contributions; P1--P3 compute the contribution for each endpoint from its
own composite. We average contributions with equal record weights within
each category and normalize only the final mean. Observed pair counts
range from 4 to 101; we retain these frequencies without reweighting.

% Generated by scripts/render_paired_ablation.py; do not edit numbers.
\begin{table}[t]
    \centering
    \footnotesize
    \setlength{\tabcolsep}{4pt}
    \begin{tabular}{lrr}
        \toprule
        Estimator & Macro AUC $\uparrow$ & Macro AP $\uparrow$ \\
        \midrule
        P0: paired, shared weights & \textbf{0.9606} & \textbf{0.7993} \\
        P1: single-sided, shared weights & 0.8818 & 0.4504 \\
        P2: raw, shared weights & 0.5102 & 0.1057 \\
        P3: single-sided, own weights & 0.8804 & 0.4494 \\
        Random directions (25 sets) & 0.5203 & 0.1136 \\
        \bottomrule
    \end{tabular}
    \caption{\textbf{Paired-intervention ablation at L18.} Frozen Qwen2.5-VL-3B-Instruct; 2,104 fitting records, 10 categories, and 5,000 COCO val2017 images. All variants share the same cosine/top-5\% readout. AUC and AP use the 0--1 scale. The random row averages metrics over 25 direction sets. The primary paired AUC difference, P0 minus P1, is $+0.0788$ (95\% CI $[+0.0717,+0.0860]$).}
    \label{tab:paired_primary}
\end{table}


\paragraph{Paired subtraction improves transfer.}
At L18, P0 reaches 0.9606 macro AUC, compared with 0.8818 for P1
(\cref{tab:paired_primary}). The paired difference is
\textbf{+7.88 AUC percentage points}, with a 95\% confidence interval
of \textbf{[+7.17, +8.60]}. We obtain percentile intervals from 2,000
image-bootstrap draws, resampling image indices jointly across categories
and variants. The interval exceeds the prespecified practical difference
of 0.5 points. Macro AP also increases, from 0.4504 to 0.7993.
P0 improves AUC over P1 for all 10 categories and at every tested layer,
with layer-wise gains of 5.05--10.77 points.

\paragraph{Single-sided reweighting does not recover the gain.}
P3 reaches 0.8804 AUC. Its difference from P1 is only \(-0.0013\) AUC
(95\% CI \([-0.0015,-0.0011]\)), whereas its deficit to P0 is 8.01
points. P2 reaches 0.5102 AUC; its difference from the measured random
control is \(-0.0101\) (95\% CI \([-0.0230,+0.0018]\)). In this
setting, the subtraction reference has a substantially larger effect than
switching between the two single-sided weighting rules.

\paragraph{Scope.}
These results support retaining paired subtraction in the tested 10-class
pipeline. They do not identify whether the gain comes from generic edit
responses, scene context, or competitor semantics. All variants retain
the historical two-category selection gates, and the intervals condition
on the fixed discovery records. Only P0 versus P1 at L18 is the primary
comparison; other comparisons are diagnostic with unadjusted intervals.
The 76-class paired-subtraction replication remains uncompleted because
the original three-image caches and complete replay recipes are unavailable.
The supplementary material provides full layer and category results,
implementation details, and replay validation.
